\documentclass[12pt]{ximera}
\typeout{Start loading xmPreamble.tex}%

% Add here extra macro's that are loaded automatically by all documents of claas 'ximera' or 'xourse' in this repo

%%
%%  Example:
%%
% \newcommand{\R}{\mathbb{R}}
\usepackage{tikz}
\usetikzlibrary{positioning, arrows.meta}
\usepackage{multicol}
\usepackage{enumitem}
\usepackage{caption}
\usepackage{amsmath}
\usepackage{amsthm}

%% NOTE: do NOT add \usepackage to individual activity .tex files.
%% When a xourse inlines an activity it gobbles \documentclass and \input, but a bare
%% \usepackage still executes -- after \begin{document} -- and errors out.
%% All shared packages and macros belong here.
%%
%% ximera.cls already loads: amsmath, amssymb, amsthm, cancel, enumitem, graphicx,
%% hyperref, listings, pgfplots, tikz, url, xcolor. No need to re-load those.

\title{In-Class: Rough Estimations}
\begin{document}
\begin{abstract}
Group work on rough estimation: budgeting a pizza party, and estimating quantities
like messages sent per year, gallons of gas to Chicago, seconds alive, daily U.S.
flights, and monthly smartphone sales --- all without a calculator.
\end{abstract}
\maketitle

\section*{Rough Estimations}

Today we will be discussing scenarios where we would like to do some rough
estimations. It is important to note a few things.

\begin{enumerate}
\item Estimations are NOT guesses. They can sometimes be educated guesses, but
  estimations, at least in this course, will always still require some sort of
  calculation.
\item Rough estimations are meant to be calculated mentally --- without calculators,
  or even pen and paper. You will still be expected to record your work and thought
  process for your instructor, but it should be work you are able to do in your head.
\item Everyone has different skill and comfort levels with mental calculation. You
  may need to round values a lot to get something you are comfortable with. That is
  completely fine, but the important thing is that in your write-up you explain
  \emph{what} numbers you rounded and \emph{why}.
\end{enumerate}

\subsection*{Pizza Party}

Let's give this a try, working with your group to determine what you would buy for a
pizza party in the following scenario.

You and your roommate are going to have some people over later, so you go to the
grocery store to grab some snacks. Everyone has agreed to pitch in \$10 to pay for
pizza and snacks for the night, and you think about 12 people are coming over. Here
are the prices of various snacks from the grocery store:

\begin{center}
\begin{tabular}{lr}
Bag of tortilla chips & \$3.99 \\
Salsa & \$1.79 \\
Bag of ``normal'' chips & \$2.99 \\
Dozen cookies from bakery & \$4.99 \\
Veggie tray & \$14.99 \\
12-pack of soda & \$5.79 \\
\end{tabular}
\end{center}

You want to get at least one of each of these items, but you're also going to order
some pizza and breadsticks. Here are the prices from your pizza place:

\begin{center}
\begin{tabular}{lr}
Cheese, medium & \$9.99 \\
Cheese, large & \$11.99 \\
Pepperoni, medium & \$11.49 \\
Pepperoni, large & \$13.49 \\
Breadsticks, 5 for & \$4.99 \\
Delivery charge & \$2.49 \\
\end{tabular}
\end{center}

Don't forget the tip!

You are a good person and do not plan to pocket any of the money your friends give
you. Plus, you and your roommate are going to pay your fair share (also \$10 each).

% TODO: in the source PDF the paragraph above is struck through by hand. If it was
% deliberately removed, the budget is 12 x $10 = $120, not 14 x $10 = $140. Confirm
% with the instructor before turning the problem below into a graded \answer{}.

\begin{problem}
Roughly how much money do you have to work with in total? State your assumption about
who is contributing.

\begin{freeResponse}
It depends who you count. If the \$10 comes only from the $12$ guests, the budget is
$12 \cdot \$10 = \$120$. If you and your roommate each pitch in as well, that is $14$
people, so $14 \cdot \$10 = \$140$.

State whichever assumption you made --- the point is that the assumption has to be
written down, because the answer changes with it.
\end{freeResponse}
\end{problem}

\begin{problem}
Decide what you're going to get. List your order, estimate the total, and explain
where you rounded.

\begin{freeResponse}
One possible order, rounding every price up to the nearest dollar as we go so the
running total stays easy:

\begin{center}
\begin{tabular}{lrr}
 & actual & rounded \\
Tortilla chips & \$3.99 & \$4 \\
Salsa & \$1.79 & \$2 \\
Normal chips & \$2.99 & \$3 \\
Cookies & \$4.99 & \$5 \\
Veggie tray & \$14.99 & \$15 \\
12-pack of soda & \$5.79 & \$6 \\
$3 \times$ large pepperoni & \$40.47 & \$42 \\
Breadsticks & \$4.99 & \$5 \\
Delivery & \$2.49 & \$3 \\
\end{tabular}
\end{center}

Rounded total: $4 + 2 + 3 + 5 + 15 + 6 + 42 + 5 + 3 = \$85$. Leaving room for a tip of
about \$15 brings us to \$100, comfortably under budget either way.

Since every price was rounded \emph{up}, this is an overestimate --- which is the safe
direction when you are checking whether you can afford something.
\end{freeResponse}
\end{problem}

\subsection*{Estimation problems}

For the following problems, do your best to determine a roughly estimated value.
Start by writing down your assumed or estimated starting values. Be sure to write
down all of your numbers and calculations, indicating anywhere you rounded. Do not
use a calculator; instead round or estimate values so you can do all calculations in
your head.

\begin{example}
Roughly estimate how many instant messages you send in a year (texts, Snapchat,
WhatsApp, etc.).

Suppose I send about $10$ texts a day. Then
\[
\frac{10 \text{ texts}}{1 \text{ day}} \cdot \frac{365 \text{ days}}{1 \text{ year}}
\approx 3650 \text{ texts in a year}.
\]
\end{example}

\begin{problem}
Roughly estimate how many instant messages \emph{you} send in a year. State your
assumed daily rate.

\begin{freeResponse}
Your daily rate will differ from the example, but the structure is the same. Suppose
you send about $40$ messages a day. Round $365$ days up to $400$ to keep the
multiplication mental:
\[
\frac{40 \text{ texts}}{1 \cancel{\text{day}}}
\cdot \frac{400 \cancel{\text{days}}}{1 \text{ year}}
= 16{,}000 \text{ texts per year}.
\]
Rounding the number of days up makes this a slight overestimate.
\end{freeResponse}
\end{problem}

\begin{problem}
Roughly estimate how much time you spend watching TV or movies in a year.

\begin{freeResponse}
Suppose about $2$ hours a day. Rounding $365$ up to $400$ days:
\[
\frac{2 \text{ hrs}}{1 \cancel{\text{day}}}
\cdot \frac{400 \cancel{\text{days}}}{1 \text{ year}}
= 800 \text{ hours per year}.
\]
It is worth converting that into a unit you have a feel for:
$800 \text{ hrs} \div 24 \approx 33$ days, so roughly a month of round-the-clock
watching per year.
\end{freeResponse}
\end{problem}

\begin{problem}
Roughly estimate how many gallons of gasoline you might use to drive from here to
Chicago, IL. Columbus, OH to Chicago, IL is about $350$ miles. Assume the vehicle
gets about $40$ mi/gal.

Roughly $\answer{9}$ gallons.

\begin{solution}
Round $350$ up to $360$ so it divides nicely by $40$:
\[
350 \cancel{\text{miles}} \cdot \frac{1 \text{ gal}}{40 \cancel{\text{miles}}}
\approx 360 \cancel{\text{miles}} \cdot \frac{1 \text{ gal}}{40 \cancel{\text{miles}}}
= \frac{36}{4} \text{ gallons} = 9 \text{ gallons}.
\]
Because we rounded the mileage \emph{up}, this is a slight overestimate.
\end{solution}
\end{problem}

\begin{problem}
Roughly estimate how much you spend on drinks (coffee, soda, tea, water, etc.) each
month.

\begin{freeResponse}
Suppose a \$5 coffee four times a week, plus about \$10 a week on other drinks --- so
roughly \$30 a week. Rounding a month up to $4$ weeks:
\[
\frac{\$30}{1 \cancel{\text{week}}} \cdot \frac{4 \cancel{\text{weeks}}}{1 \text{ month}}
= \$120 \text{ per month}.
\]
A month is really closer to $4.3$ weeks, so this is a slight underestimate --- call it
\$120 to \$130.
\end{freeResponse}
\end{problem}

\begin{problem}
Roughly estimate how many seconds you've been alive. Let's say you are $20$ years old.

Roughly $\answer{600000000}$ seconds.

\begin{hint}
First write down the structure, then round at each step:
\[
20 \text{ years} \cdot \frac{365 \text{ days}}{1 \text{ year}}
\cdot \frac{24 \text{ hrs}}{1 \text{ day}}
\cdot \frac{60 \text{ min}}{1 \text{ hr}}
\cdot \frac{60 \text{ sec}}{1 \text{ min}}.
\]
\end{hint}

\begin{solution}
Rounding as we go:
\[
\begin{aligned}
20 \cancel{\text{years}} \cdot \frac{400 \text{ days}}{1 \cancel{\text{year}}}
  &= 8000 \text{ days} \\
8000 \cancel{\text{days}} \cdot \frac{25 \text{ hrs}}{1 \cancel{\text{day}}}
  &= 200{,}000 \text{ hrs} \\
200{,}000 \cancel{\text{hrs}} \cdot \frac{60 \text{ min}}{1 \cancel{\text{hr}}}
  &= 12{,}000{,}000 \text{ min} \approx 10{,}000{,}000 \text{ min} \\
10{,}000{,}000 \cancel{\text{min}} \cdot \frac{60 \text{ sec}}{1 \cancel{\text{min}}}
  &= 600{,}000{,}000 \text{ sec}.
\end{aligned}
\]
So a $20$ year old has been alive for roughly $600$ million seconds.

Note the deliberate round-down from $12$ million to $10$ million minutes: it keeps the
last multiplication easy. Rounding both up and down along the way means we can no
longer say for certain whether the final estimate is high or low.
\end{solution}
\end{problem}

\begin{problem}
Roughly estimate how many flights you think there are in the U.S. daily. Helpful
piece of information: CMH (Columbus airport) averages about $140$ flights daily.

\begin{freeResponse}
Start from CMH and scale up. Columbus is a mid-sized airport, so use it as a typical
one and estimate how many airports of that size the country has. Suppose about $500$
airports with meaningful commercial traffic. Round $140$ down to $100$ to keep it
mental:
\[
\frac{100 \text{ flights}}{1 \cancel{\text{airport}}} \cdot 500 \cancel{\text{airports}}
= 50{,}000 \text{ flights a day}.
\]
Notice how much rests on the assumed number of airports --- that is the number worth
arguing about with your group, not the arithmetic. Any answer in the tens of
thousands is defensible.
\end{freeResponse}
\end{problem}

\begin{problem}
Roughly estimate how many smartphones you think are sold in the U.S. each month.
There are about $325$ million people in the U.S.

\begin{freeResponse}
Round $325$ million up to $300$ million people --- actually down, which keeps the
division easy. Not everyone owns a smartphone, but most do, so assume roughly $300$
million phones in use. If people replace a phone about every $3$ years, that is every
$36$ months, which we round to $30$:
\[
\frac{300{,}000{,}000 \text{ phones}}{30 \text{ months}}
= 10{,}000{,}000 \text{ phones per month}.
\]
So roughly $10$ million a month. The replacement interval is the assumption doing the
real work here --- if people keep phones $4$ years instead of $3$, the estimate drops
by a quarter.
\end{freeResponse}
\end{problem}

\begin{problem}
Look back at the ``seconds alive'' and ``smartphones sold'' questions.

\begin{prompt}
\textbf{(a)} For each, determine a number that would definitely be way too low for
even a rough estimate, but that still requires some calculation. Explain why you
think this would be unreasonably low. Still no calculator --- use mental arithmetic.
\begin{freeResponse}
\textbf{Seconds alive.} One million seconds sounds enormous, but check it:
$1{,}000{,}000 \div 60 \approx 17{,}000$ minutes, $\div 60 \approx 280$ hours,
$\div 24 \approx 12$ days. A $20$ year old has been alive far longer than $12$ days,
so a million is much too low.

\textbf{Smartphones per month.} One thousand a month would mean roughly one phone for
every $325{,}000$ people --- about one phone per month for a city the size of
Cincinnati. Far too low.

Note that in both cases the number only looks unreasonable \emph{after} you convert it
into units you have intuition for. That conversion is the calculation.
\end{freeResponse}
\end{prompt}

\begin{prompt}
\textbf{(b)} For each, determine a number that would definitely be way too high for
even a rough estimate. Explain why you think this would be unreasonably high.
\begin{freeResponse}
\textbf{Seconds alive.} One hundred billion seconds: a year is roughly
$400 \cdot 25 \cdot 60 \cdot 60 \approx 36$ million seconds, so
$100{,}000{,}000{,}000 \div 36{,}000{,}000 \approx 3000$ years. Nobody is $3000$ years
old, so this is far too high.

\textbf{Smartphones per month.} One hundred million a month would be $1.2$ billion a
year, in a country of $325$ million people --- close to four new phones per person per
year. Far too high.
\end{freeResponse}
\end{prompt}
\end{problem}

\end{document}
